\section{Discussion and Policy Implications}

\subsection{The Time Horizon Imperative}

The 100-year optimization reveals a stark departure from shorter-horizon analyses: zero deforestation emerges as economically optimal when society accounts for century-scale forest values. Current practice—58,000 ha/year deforestation—destroys economic value equivalent to \$19.7 billion in present terms, representing the welfare cost of policy inaction. This finding upends conventional wisdom that economic development requires tolerating some forest loss.

Why does extended time horizon shift the prescription so dramatically? Three mechanisms operate. First, terminal value—forest stock in year 100 valued at \$14,350/ha through cumulative ecosystem services—becomes large enough to dominate near-term agricultural returns. Second, carbon sequestration compounds over a century; each hectare reforested in 2024 sequesters 150 Mg CO$_2$e by 2124, generating \$3,000 in present-value climate benefits. Third, the logistic growth function implies that forests near carrying capacity regenerate slowly, making early conservation critical to achieving long-run equilibrium.

Yet farmers operate under dramatically shorter planning horizons. Credit constraints, insecure tenure, and immediate subsistence needs generate implicit discount rates often exceeding 20\% annually—sevenfold higher than the 3\% social rate we employ. This temporal wedge—not ignorance or irrationality—drives the private-social divergence. Policy must bridge this gap through mechanisms that transfer long-run social values into near-term private incentives.

\subsection{Regulatory Approaches: Enforcement Constraints}

Ghana maintains a legal framework restricting forest clearing. The 1927 Forest Ordinance designates 266 forest reserves covering 2.1 million hectares as protected; the 1974 Control and Prevention of Bushfires Law prohibits unauthorized burning; the 2016 Logging and Timber Harvesting Regulations require permits for tree felling. On paper, Ghana's forest law appears comprehensive.

Enforcement collapses in practice. The Forestry Commission employs approximately 1,200 field staff to monitor 9.5 million hectares—one officer per 7,900 hectares. For comparison, effective protected area management typically requires one ranger per 1,000-2,000 hectares \citep{jachmann2008monitoring}. Ghana's staffing ratio falls short by a factor of five.

Budget constraints worsen the problem. Real expenditure on forest protection declined 34\% between 2015 and 2024 after accounting for inflation. Field stations lack vehicles; officers patrol on foot or bicycle, covering limited territory. Remote sensing identifies illegal clearing rapidly, but ground response lags by months—after which farmers have planted cocoa and enforcement becomes politically contentious.

Even when violations are detected, prosecution fails. Between 2018 and 2023, the Forestry Commission documented 3,847 cases of illegal clearing in forest reserves. Prosecutions numbered 127 (3.3\%), resulting in 41 convictions. Median fine: 500 cedis (\$35), far below the agricultural revenue from cleared land. When expected penalties fall below expected gains, deterrence evaporates.

Could enforcement be strengthened? Yes, but at substantial cost. Doubling field staff to minimally adequate levels requires an additional \$18 million annually. Equipping stations with vehicles and establishing effective patrol networks adds \$12 million capital investment. Total: \$30 million per year. Is this feasible? Ghana's 2024 Forestry Commission budget totaled \$22 million—the required increase exceeds current spending by 136\%.

Alternative enforcement mechanisms warrant exploration. Community forest management programs transfer monitoring responsibility to local villages, who have stronger incentives to prevent outsider encroachment. \citet{bowler2012community} finds that community-managed forests experience 29\% less deforestation than government reserves, attributing this to superior local knowledge and reduced enforcement costs. Ghana has piloted such programs in 38 communities; scaling nationally could improve outcomes at lower fiscal cost than expanding central enforcement.

\subsection{Carbon Pricing: Threshold Effects and Market Access}

Our sensitivity analysis identifies \$15/Mg as the critical threshold where optimal deforestation drops to zero under 100-year horizons—substantially lower than the \$28/Mg threshold emerging from 20-year models. This difference illustrates how terminal value magnifies carbon pricing effectiveness across extended periods. Current voluntary market prices (\$12-\$20/Mg) sit near this threshold, suggesting modest price increases could trigger discontinuous behavioral shifts.

Yet access barriers prevent most Ghanaian farmers from capturing carbon revenues. Transaction costs—baseline assessment (\$50,000-\$100,000), monitoring (\$20,000/year), verification (\$30,000 per cycle)—make projects below 10,000 ha economically nonviable. Ghana's average forest parcel spans 2.3 ha, requiring aggregation of 4,300 farmers to reach minimum scale. Coordination failures plague such efforts: liability allocation when members defect, equitable revenue distribution, and leadership accountability all generate bargaining costs that consume potential rents.

Simplified protocols might help. The World Bank's Climate Warehouse initiative cuts verification costs to \$15,000 per project through satellite monitoring, dropping viable scale to 2,000 ha. But additionality concerns intensify under simplified methods—paying farmers not to clear forests they would preserve anyway wastes resources. The optimal protocol balances verification rigor against transaction cost minimization, a trade-off requiring empirical calibration through pilot programs.

\subsection{Agricultural Intensification: Relieving Land Pressure}

Deforestation pressure traces fundamentally to agricultural land demand. Ghana's cocoa output expanded from 340,000 to 850,000 metric tons (2000-2024); population grew from 19 to 34 million. Absent yield improvements, this growth trajectory requires proportional land expansion. Yet yields average 500 kg/ha against a technical frontier of 2,200 kg/ha—a fourfold gap indicating vast intensification potential.

Closing half this gap would allow current output on one-quarter the land area, freeing 600,000 ha for reforestation. Our calculations suggest investing \$500 million in intensification programs—financing replanting, subsidizing inputs, expanding extension services—eliminates the need to clear 225,000 ha over the century. At \$410/ha in foregone environmental damages, this prevents \$92 million in harm—an 18\% benefit-cost ratio excluding increased cocoa revenues from higher yields.

Three constraints bind intensification. Aging tree stock (67\% exceed 25 years) operates past peak productivity; replanting requires 3-4 year income sacrifice that credit-constrained farmers cannot bear. Limited input adoption (23\% use fertilizer, 12\% improved varieties) reflects both capital constraints and knowledge gaps—extension services reach only 18\% of farmers annually. Cocoa swollen shoot virus afflicting 250,000 ha creates incentives to abandon diseased plantations and expand elsewhere rather than rehabilitate affected areas.

Policy levers addressing each constraint exist: credit programs financing replanting with income support during unproductive periods; mobile extension services diffusing best practices at scale; subsidized CSSV treatment with strict quarantine enforcement. The obstacle is not technical but institutional—agricultural ministries prioritize food security and export earnings, environment ministries lack budgets to influence farming practices. Effective intensification requires coordination across ministries whose incentives diverge, a governance challenge as binding as market failures.

\subsection{Reforestation at Scale: Institutional Bottlenecks}

Optimal policy prescribes aggressive reforestation—264,000 ha/year initially, declining as forest stock approaches carrying capacity. Ghana's actual rates (11,500 ha/year, 2015-2024) fall 96% short, revealing stark implementation gaps. The National Forest Plantation Development Programme targets 250,000 ha by 2040; through 2024, it established 41,000 ha—16% of interim targets.

Multiple bottlenecks constrain scaling. Land tenure uncertainty affects 72\% of rural areas lacking formal title, discouraging investment in tree planting with 15-20 year return horizons. Seedling quality limits effective supply: government nurseries produce 15 million stems annually, but 38% survival rates reduce usable stock to 5.7 million—sufficient for only 3,800 ha at standard densities.

Financing mechanisms remain underdeveloped. Establishment costs of \$2,000/ha require upfront capital farmers lack; returns arrive slowly while maintenance demands continuous effort. Payment for ecosystem services could bridge this temporal mismatch—offering annual payments during establishment—but Ghana's PES schemes remain pilots covering under 5,000 ha nationally.

Natural regeneration offers lower-cost alternatives. Our data document 229,400 ha regenerating naturally (2000-2020) on abandoned agricultural land at zero planting cost. Policies facilitating longer fallow rotations—perhaps through crop insurance maintaining farmer income during non-cultivation—could amplify natural processes. Yet secondary forests require 40-60 years to match primary forest biomass and rarely achieve comparable biodiversity. Strategic active reforestation targeting corridor creation, riparian restoration, and slope stabilization delivers ecosystem benefits natural regeneration cannot replicate.

\subsection{Threshold Dynamics and Irreversibility}

Figure \ref{fig:comparison} illustrates a critical insight: business-as-usual crosses ecological thresholds that optimal policy avoids entirely. The connectivity threshold (2 million ha) marks the point where forest fragmentation disrupts seed dispersal, pollinator movements, and microclimate regulation. BAU trajectories cross this threshold around 2072; once crossed, restoration costs spike and success rates plummet as ecological processes break down.

This irreversibility generates option value for conservation—maintaining forest above thresholds preserves future management flexibility. Our model captures this imperfectly through the logistic growth function, but reality features sharper discontinuities. Below 2 million ha, Ghana's forests fragment into isolated patches unable to maintain viable populations of large-seeded tree species dependent on dispersal by now-extirpated megafauna. Reforestation in such landscapes requires artificial seeding at substantially higher cost.

The NDC threshold (5.5 million ha) represents a political commitment under the Paris Agreement. BAU violates this by 2051, triggering potential diplomatic and financial penalties as Ghana loses eligibility for climate finance conditional on emission reductions. The optimal path maintains stock above 6 million ha throughout, ensuring compliance while generating sequestration credits eligible for international carbon markets.

These dynamics underscore time-inconsistency problems in forest policy. Optimal long-run management prescribes immediate action—crossing thresholds imposes irreversible costs—but political incentives favor delay. Overcoming this requires credible commitment mechanisms: multi-decade budget allocations, legally binding conservation targets, and institutional designs that insulate forest policy from short-term political cycles.

\subsection{Political Economy and Implementation Feasibility}

Why have economically justified interventions not materialized? Political economy constraints bind tighter than resource constraints. Cocoa farmers—800,000 strong—wield electoral influence; forest benefits accrue diffusely across constituencies. Politicians maximizing near-term support rationally prioritize farmer interests over long-run environmental outcomes.

Corruption exacerbates enforcement failures. Forestry officers earning modest salaries face tempting bribes to ignore illegal clearing. Evidence from timber inspection checkpoints documents that 67\% accept payments to pass illegally logged wood; comparable clearing violations likely follow similar patterns. Civil service reform—higher salaries, stronger oversight, performance incentives—is necessary but politically difficult.

International finance offers potential catalysts. Norway's climate fund has disbursed \$1.1 billion to forest countries since 2008, contingent on verified emission reductions. Ghana has not accessed these funds at scale because deforestation increased rather than decreased since joining REDD+. Result-based payments create powerful incentives, but require achieving results first—a chicken-egg problem for countries lacking implementation capacity.

Corporate commitments provide another leverage point. Major cocoa traders pledged zero-deforestation supply chains, creating potential conservation allies. Yet commitments contain loopholes—"zero net deforestation" permits clearing offset by reforestation elsewhere—and traceability systems remain incomplete. Only 43\% of Ghana's cocoa traces to farm-level origins; the remainder moves through opaque aggregation chains where deforestation-linked production mingles freely with compliant cocoa.

\subsection{Integrated Policy Framework}

No single intervention suffices; optimal outcomes require coordinated action across multiple domains. Based on our analysis, we propose a sequenced implementation strategy:

\textbf{Immediate priorities (Years 1-3)}: Strengthen enforcement in critical biodiversity zones through community forest management, starting with 50,000 ha in Ankasa and Kakum Conservation Areas where local monitoring costs one-fifth of central enforcement. Launch agricultural intensification pilots in high-deforestation districts, targeting 20,000 farmers with replanting support and input subsidies.

\textbf{Medium-term buildout (Years 4-7)}: Scale successful intensification programs nationally, aiming for 200,000 farmers adopting improved practices. Develop simplified carbon credit protocols and aggregate smallholders into 10-15 projects qualifying for international carbon finance. Establish 100,000 ha of strategic reforestation in corridors connecting isolated forest fragments.

\textbf{Long-term institutionalization (Years 8-20)}: Implement spatially differentiated land-use zoning with transparent rules and sustained enforcement capacity. Transition carbon programs from donor dependence to self-financing through credit revenues. Achieve forest stock stabilization at 8.5-9.0 million ha through balanced natural regeneration and minimal active management.

Total costs estimate \$1.2 billion over 100 years (\$60 million annually), equivalent to 0.08\% of Ghana's GDP. Benefits—avoided carbon emissions (\$80 billion valued at \$20/Mg over the century), watershed protection, biodiversity conservation—generate benefit-cost ratios exceeding 7:1 even under conservative assumptions. The question is not whether Ghana can afford conservation; it is whether political systems can organize collective action to implement known solutions.
